% ==============================================================================
% CHƯƠNG 1: TỔNG QUAN ĐỀ TÀI
% ==============================================================================

\chapter{Tổng quan đề tài}
\label{chap:tong-quan-de-tai}

\section{Bối cảnh và lý do lựa chọn đề tài}
\label{sec:boi-canh-ly-do}

Trong hai thập kỷ qua, khối lượng tài liệu khoa học số tăng nhanh vượt quá khả năng đọc và tổng hợp thủ công. Các công cụ tìm kiếm dựa trên từ khóa truyền thống vẫn tồn tại ba hạn chế lớn: khó nắm bắt ý nghĩa ngữ nghĩa, buộc người dùng tự đánh giá độ liên quan của từng tài liệu và chưa hỗ trợ hỏi đáp tổng hợp kèm dẫn nguồn kiểm chứng được.

Sự phát triển của mô hình ngôn ngữ lớn cùng kiến trúc Retrieval-Augmented Generation (RAG) mở ra hướng tiếp cận mới bằng cách truy hồi các đoạn văn bản liên quan trước khi sinh câu trả lời. Tuy nhiên, ứng dụng RAG cho bài báo khoa học vẫn đầy thách thức do nội dung chuyên sâu, cấu trúc dài và yêu cầu khắt khe về độ chính xác cũng như khả năng truy vết nguồn.

Từ thực tế đó, đề tài xây dựng hệ thống RAG hỗ trợ tìm kiếm ngữ nghĩa và tra cứu tương tác đối với bài báo khoa học, sử dụng thống nhất TypeScript cho cả phía máy chủ và phía người dùng. Việc dùng chung một ngôn ngữ giúp giảm chi phí phối hợp, chia sẻ kiểu dữ liệu và đồng nhất cơ chế kiểm tra đầu vào trong toàn hệ thống.

\section{Giới thiệu bài toán và định hướng giải pháp}
\label{sec:gioi-thieu-bai-toan}

Bài toán đặt ra là xây dựng hệ thống cho phép người dùng truy vấn kho tài liệu khoa học bằng ngôn ngữ tự nhiên, truy hồi các đoạn văn bản liên quan và dùng mô hình ngôn ngữ lớn tổng hợp câu trả lời kèm dẫn nguồn kiểm chứng được. Giải pháp được tổ chức thành bốn lớp chức năng:

\begin{itemize}
  \item[--] \textbf{Lớp dữ liệu và lập chỉ mục}: Thu thập, làm sạch, khử trùng lặp và chuẩn hóa bài báo từ nguồn công khai như arXiv; phân đoạn văn bản, gắn siêu dữ liệu (metadata) và chuyển thành vector nhúng để phục vụ tìm kiếm ngữ nghĩa.
  \item[--] \textbf{Lớp truy hồi}: Chuyển câu hỏi thành vector, tìm kiếm các đoạn tương đồng trong cơ sở dữ liệu vector và xếp hạng lại bằng các tín hiệu bổ sung (như độ uy tín tác giả từ đồ thị trích dẫn) để chọn ngữ cảnh phù hợp.
  \item[--] \textbf{Lớp sinh câu trả lời}: Cung cấp ngữ cảnh truy hồi cho mô hình ngôn ngữ lớn theo mẫu câu lệnh có kiểm soát, phát câu trả lời theo thời gian thực kèm định danh nguồn và kiểm tra tính hợp lệ của trích dẫn.
  \item[--] \textbf{Lớp giao diện và trải nghiệm người dùng (UI/UX)}: Kết hợp không gian trực quan hóa cây tri thức ba chiều bằng WebGL với khung trò chuyện học thuật; lịch sử hội thoại được lưu cục bộ theo cơ chế ưu tiên ngoại tuyến (offline-first) và tự động đồng bộ khi có kết nối.
\end{itemize}

\section{Mục tiêu xây dựng hệ thống}
\label{sec:muc-tieu}

Mục tiêu tổng quát của đề tài là xây dựng một hệ thống tra cứu tài liệu khoa học hoàn chỉnh, hỗ trợ hỏi đáp bằng ngôn ngữ tự nhiên, khám phá trực quan và đảm bảo khả năng truy vết nguồn cho mọi câu trả lời. Các mục tiêu cụ thể bao gồm:

\begin{itemize}
  \item[--] \textbf{Thu thập và lập chỉ mục ngữ nghĩa}: Xây dựng quy trình thu thập, khử trùng lặp và chuẩn hóa kho bài báo kèm siêu dữ liệu đầy đủ; thiết kế chiến lược phân đoạn, tạo vector nhúng và lưu trữ tối ưu cho truy vấn ngữ nghĩa.
  \item[--] \textbf{Truy hồi nhiều tầng}: Kết hợp tìm kiếm vector với xếp hạng lại dựa trên tín hiệu đồ thị trích dẫn nhằm nâng cao chất lượng tập ứng viên đưa vào ngữ cảnh sinh.
  \item[--] \textbf{Sinh câu trả lời cục bộ có kiểm chứng}: Vận hành mô hình ngôn ngữ lớn hoàn toàn trong hạ tầng cục bộ, hỗ trợ phát trực tiếp theo luồng và dẫn nguồn kiểm chứng được, bảo đảm dữ liệu người dùng không rời khỏi hệ thống.
  \item[--] \textbf{Giao diện học thuật và lưu trữ ngoại tuyến}: Xây dựng không gian khám phá ba chiều cùng khung trò chuyện hiển thị công thức toán học và trích dẫn; lưu lịch sử hội thoại ưu tiên ngoại tuyến trên trình duyệt và đồng bộ hai chiều với máy chủ.
  \item[--] \textbf{Xác thực, phân quyền và kiểm toán}: Triển khai băm mật khẩu theo Argon2id, cấp và xoay vòng cặp JWT truy cập/làm mới theo họ token (token family), kèm nhật ký quản trị chỉ bổ sung (append-only).
  \item[--] \textbf{Kiểm thử và đánh giá toàn diện}: Xây dựng bộ kiểm thử đơn vị, tích hợp, đầu cuối cùng quy trình đánh giá chất lượng truy hồi, chất lượng sinh nội dung và hiệu năng dưới tải.
\end{itemize}

\section{Phạm vi và giới hạn của đề tài}
\label{sec:pham-vi-gioi-han}

Phạm vi thực hiện của đề tài được xác định theo ba trục chính:

\begin{itemize}
  \item[--] \textbf{Nghiệp vụ}: Hỗ trợ đăng ký, đăng nhập, tìm kiếm ngữ nghĩa và khám phá trực quan, quản lý nhiều phiên hội thoại có dẫn nguồn, đồng bộ lịch sử giữa trình duyệt và máy chủ, cùng bảng quản trị người dùng, phiên đăng nhập và nhật ký kiểm toán.
  \item[--] \textbf{Kiến trúc}: Tổ chức gồm ứng dụng máy chủ, ứng dụng web trình duyệt và tầng dữ liệu tách biệt (cơ sở dữ liệu hướng tài liệu, cơ sở dữ liệu vector và mô hình ngôn ngữ cục bộ); vận hành tự chủ không phụ thuộc dịch vụ bên thứ ba và triển khai được trên một máy chủ đơn.
  \item[--] \textbf{Dữ liệu}: Tập trung vào kho bài báo bản tiền ấn (preprint) công khai từ arXiv thuộc các lĩnh vực trí tuệ nhân tạo, học máy và ngành liên quan, được thu thập tăng dần và phiên bản hóa.
\end{itemize}

Các giới hạn của đề tài bao gồm:

\begin{itemize}
  \item[--] \textbf{Nguồn dữ liệu và ngôn ngữ}: Giới hạn ở tài liệu preprint công khai bằng tiếng Anh, không bao gồm tạp chí trả phí hay nguồn độc quyền (dù giao diện sẵn sàng mở rộng đa ngôn ngữ).
  \item[--] \textbf{Định dạng trích dẫn}: Dẫn nguồn theo mã định danh tài liệu để đối chiếu bản gốc, không thay thế bộ chuẩn hóa định dạng trích dẫn học thuật đầy đủ.
  \item[--] \textbf{Hiệu năng và phần cứng}: Ưu tiên tính chính xác và mức độ có căn cứ hơn tốc độ sinh tối đa; thông lượng phụ thuộc trực tiếp vào tài nguyên phần cứng cục bộ.
  \item[--] \textbf{Phạm vi học thuật và cộng tác}: Không thay thế vai trò kiểm định hay phản biện khoa học của con người, và chưa hỗ trợ chia sẻ phiên, cộng tác nhóm hay gợi ý cá nhân hóa.
\end{itemize}

\section{Đối tượng sử dụng và các tác nhân liên quan}
\label{sec:doi-tuong-tac-nhan}

Hệ thống hướng đến bốn nhóm đối tượng sử dụng chính:

\begin{itemize}
  \item[--] \textbf{Nhà nghiên cứu và nghiên cứu sinh}: Cần rà soát nhanh công trình liên quan, nắm ý chính bài báo và truy vết nguồn đáng tin cậy trước khi trích dẫn.
  \item[--] \textbf{Sinh viên khoa học -- kỹ thuật}: Cần giải thích khái niệm dễ hiểu, tìm tài liệu nền tảng và đọc công thức toán học trực quan.
  \item[--] \textbf{Kỹ sư và chuyên gia phân tích}: Cần lọc nhanh tài liệu theo chủ đề, so sánh cách tiếp cận và quan sát xu hướng qua không gian tri thức.
  \item[--] \textbf{Tác nhân tự động}: Tích hợp truy xuất và tóm tắt tài liệu theo lập trình thông qua giao diện lập trình ứng dụng (API).
\end{itemize}

Về phía hệ thống, có bốn nhóm tác nhân gồm: người dùng chưa đăng nhập (chỉ khám phá công khai, không lưu đồng bộ), người dùng đã đăng nhập (sử dụng đầy đủ chức năng và quản lý phiên cá nhân), quản trị viên (quản lý tài khoản, thu hồi phiên và tra cứu nhật ký) và các thành phần nội bộ (bộ thu thập dữ liệu, dịch vụ nhúng, tiến trình suy luận). Về ranh giới tin cậy, mọi yêu cầu từ bên ngoài đều qua lớp xác thực và phân quyền; mọi thao tác quản trị được ghi vào nhật ký chỉ bổ sung (append-only), còn các thành phần nội bộ chỉ được cấp quyền tối thiểu theo chế độ chỉ đọc hoặc qua hàng đợi giới hạn.


\section{Tổng quan công nghệ và công cụ sử dụng}
\label{sec:tong-quan-cong-nghe}

Bảng~\ref{tab:tong-quan-cong-nghe} trình bày tổng quan các công nghệ và công cụ được sử dụng trong
hệ thống, kèm theo vai trò và lý do lựa chọn.

\begin{longtable}[]{@{}
  >{\raggedright\arraybackslash}p{0.18\dimexpr\linewidth-8\tabcolsep\relax}
  >{\raggedright\arraybackslash}p{0.19\dimexpr\linewidth-8\tabcolsep\relax}
  >{\raggedright\arraybackslash}p{0.10\dimexpr\linewidth-8\tabcolsep\relax}
  >{\raggedright\arraybackslash}p{0.26\dimexpr\linewidth-8\tabcolsep\relax}
  >{\raggedright\arraybackslash}p{0.27\dimexpr\linewidth-8\tabcolsep\relax}@{}}
\caption{Tổng quan công nghệ và công cụ sử dụng}\label{tab:tong-quan-cong-nghe}\\
\toprule\noalign{}
\begin{minipage}[b]{\linewidth}\raggedright
\textbf{Nhóm}
\end{minipage} & \begin{minipage}[b]{\linewidth}\raggedright
\textbf{Công nghệ hoặc công cụ}
\end{minipage} & \begin{minipage}[b]{\linewidth}\raggedright
\textbf{Phiên bản}
\end{minipage} & \begin{minipage}[b]{\linewidth}\raggedright
\textbf{Vai trò trong hệ thống}
\end{minipage} & \begin{minipage}[b]{\linewidth}\raggedright
\textbf{Lý do lựa chọn chính}
\end{minipage} \\
\midrule\noalign{}
\endfirsthead
\caption[]{(tiếp theo)}\\
\toprule\noalign{}
\begin{minipage}[b]{\linewidth}\raggedright
\textbf{Nhóm}
\end{minipage} & \begin{minipage}[b]{\linewidth}\raggedright
\textbf{Công nghệ hoặc công cụ}
\end{minipage} & \begin{minipage}[b]{\linewidth}\raggedright
\textbf{Phiên bản}
\end{minipage} & \begin{minipage}[b]{\linewidth}\raggedright
\textbf{Vai trò trong hệ thống}
\end{minipage} & \begin{minipage}[b]{\linewidth}\raggedright
\textbf{Lý do lựa chọn chính}
\end{minipage} \\
\midrule\noalign{}
\endhead
\bottomrule\noalign{}
\endlastfoot
Ngôn ngữ lập trình & TypeScript & 5.7 & Toàn bộ tầng máy chủ và tầng
người dùng & Đồng nhất ngôn ngữ trong toàn hệ thống, giảm chi phí phối
hợp, kiểu tĩnh giúp kiểm soát chất lượng mã nguồn quy mô lớn \\
Nền tảng máy chủ & NestJS & 11 & Định khung nghiệp vụ, tiêm phụ thuộc,
kiểm định dữ liệu đầu vào & Tổ chức mã nguồn theo mô đun với ranh giới
rõ ràng, chuẩn hóa việc xử lý yêu cầu và phản hồi lỗi \\
Tài liệu giao diện lập trình & OpenAPI và Swagger UI & 3.0 & Sinh tự
động tài liệu và mô tả hợp đồng giao diện & Tạo nguồn tham chiếu máy đọc
cho kiểm thử hợp đồng, giảm sai lệch giữa mã nguồn và tài liệu \\
Cơ sở dữ liệu hướng tài liệu & MongoDB & 7.0 & Lưu trữ người dùng,
phiên trò chuyện, tin nhắn và nhật ký quản trị & Mô hình tài liệu linh
hoạt phù hợp dữ liệu văn bản không đồng nhất, hỗ trợ lọc theo phạm vi và
sắp xếp theo trục thời gian \\
Lớp truy xuất và mô hình dữ liệu & Mongoose & 8.x & Ánh xạ lớp và tài
liệu, định nghĩa lược đồ và chỉ mục & Kiểm soát kiểu dữ liệu ở tầng ứng
dụng, giảm lỗi hình thức khi truy cập cơ sở dữ liệu \\
Cơ sở dữ liệu vector & Qdrant & 1.x & Lưu trữ vector đoạn văn bản và phục
vụ truy hồi gần đúng theo khoảng cách cosine & Hỗ trợ chỉ mục HNSW, độ trễ thấp, khoảng cách cosine khớp trực tiếp với cách
huấn luyện mô hình nhúng \\
Mô hình nhúng văn bản & Nomic Embed & 1.5 & Biến đổi câu hỏi và đoạn văn
bản thành vector cố định chiều & Hỗ trợ tiền tố tác vụ (task prefix) phân biệt câu hỏi và
tài liệu, chuẩn hóa vector nhất quán với độ đo cosine \\
Thư viện suy luận nhúng & Transformers.js & 3.x & Chạy mô hình nhúng
trong tiến trình Node.js & Thực thi mô hình dạng ONNX phía máy chủ,
không cần dịch vụ suy luận từ xa \\
Mô hình sinh nội dung & Qwen 2.5 Instruct, định dạng GGUF lượng tử hóa
Q4 & 7 tỷ tham số & Sinh câu trả lời trong ngữ cảnh truy hồi & Cân bằng
giữa chất lượng sinh và yêu cầu bộ nhớ, tương thích với mô hình ngôn ngữ
được tinh chỉnh theo hướng dẫn (instruction-tuned) \\
Thư viện chạy mô hình cục bộ & node-llama-cpp & 3.x & Nạp trọng số, chuyển
tầng suy luận sang GPU (GPU offloading), phát trực tiếp từng mảnh kết quả & Hỗ trợ
CUDA và Metal, cơ chế hàng đợi giới hạn, cho phép chạy mô hình hoàn toàn
trong hạ tầng cục bộ \\
Thư viện giao diện người dùng & React & 19 & Dựng giao diện trình duyệt
tài liệu, khung trò chuyện và bảng quản trị & Mô hình thành phần phù hợp
giao diện nhiều trạng thái, hệ sinh thái thư viện phong phú \\
Công cụ đóng gói và phát triển & Vite & 6 & Đóng gói ứng dụng và phục vụ
môi trường phát triển & Khởi động nhanh, tăng tốc cập nhật nóng (HMR), cấu
hình đơn giản \\
Thư viện đồ họa ba chiều & three.js & r160 & Dựng cảnh không gian tài
liệu trong môi trường WebGL & Nền tảng WebGL phổ biến, hiệu năng cao, hỗ
trợ đầy đủ các thao tác hình học và ánh sáng \\
Lớp liên kết đồ họa với giao diện & React Three Fiber & 8.x & Khai báo
cảnh ba chiều bằng mô hình thành phần React & Đồng nhất ngôn ngữ khai
báo giữa phần hai chiều và phần ba chiều, quản lý vòng đời tài nguyên đồ
họa \\
Quản lý trạng thái phía người dùng & Zustand & 4.x & Lưu trạng thái trò
chuyện, vòng đời phiên và dữ liệu đồ họa & Cập nhật cục bộ không gây
hiệu ứng phụ, tránh tái kết xuất toàn bộ cây thành phần khi dữ liệu đang
phát trực tiếp \\
Lưu trữ cục bộ trình duyệt & Dexie.js trên nền tảng IndexedDB & 4.x &
Lưu trữ hội thoại và siêu dữ liệu cục bộ theo cơ chế ưu tiên ngoại tuyến
& API bất đồng bộ dựa trên Promise (Promise-based API) cho cơ sở dữ liệu cục bộ, khai báo lược
đồ và chỉ mục thuận tiện \\
Khung giao diện CSS & Tailwind CSS & 4.x & Dựng hệ thống giao diện và
các thành phần trực quan & Hệ thống tiện ích nhất quán, hỗ trợ thiết kế
theo hướng dẫn hình ảnh và giảm mã CSS thủ công \\
Hiển thị công thức toán học & KaTeX & 0.16.x & Định dạng phương trình
trong nội dung trả lời và trích dẫn & Kết xuất nhanh, đồng bộ, không phụ
thuộc phông chữ động \\
Xác thực và phân quyền & JSON Web Token, Argon2id & RFC 7519, RFC 9106 &
Quản lý đăng nhập, cấp token truy cập, xoay vòng token làm mới & Xoay
vòng token theo họ phát hiện việc tái sử dụng, Argon2id tăng đáng kể chi
phí dò mật khẩu \\
Định dạng lỗi giao diện & Problem Details for HTTP APIs & RFC 7807 &
Chuẩn hóa cấu trúc thông báo lỗi trả về cho người dùng & Thông báo lỗi
nhất quán và có thể máy đọc, thuận tiện cho cả giao diện lập trình và
người dùng cuối \\
Kiểm thử đơn vị phía máy chủ & Jest & 29.x & Kiểm thử đơn vị và tích hợp
các mô đun máy chủ & Hệ sinh thái mock sẵn có, chạy song song nhanh,
tích hợp sẵn với quy trình kiểm thử của NestJS \\
Kiểm thử đơn vị phía người dùng & Vitest & 2.x & Kiểm thử hàm, hook và
thành phần phía người dùng & Tương thích với hệ sinh thái Vite, dùng
chung đường dẫn cấu hình với công cụ đóng gói \\
Kiểm thử đầu cuối & Playwright & 1.x & Kiểm thử luồng người dùng trên
trình duyệt thật & Kiểm thử ổn định với nhiều trình duyệt, hỗ trợ chạy
song song và theo dõi hiệu năng \\
Kiểm thử tải & k6 & 1.x & Mô phỏng tải nhiều cấp độ lên các endpoint của
hệ thống & Mô hình kịch bản tải dạng mã, báo cáo độ trễ theo phân vị và
tỷ lệ lỗi \\
\end{longtable}
